This study ensured the confidentiality of all associated entities by strictly adhering to ethical data use standards and the provisions of the Data Privacy Act of 2012 (Republic Act No. 10173). Consistent with the Act’s mandate on the protection of personal information, no private individual or household data was collected or utilized. The model was predicated upon synthesized socio-demographic data derived exclusively from public aggregates, such as official figures from the Philippine Statistics Authority \citep{PSA2024}, and anonymized operational details obtained from LGU and barangay officials. This approach allowed the model to capture the necessary geographic and demographic complexity \citep{Jimenez2025} while fully protecting the privacy of citizens and local officials.

This research also acknowledged its inherent limitations. The computational model was an abstraction of reality \citep{Brugiere2022} and, as such, functioned only as a simplification. This research also acknowledged its inherent limitations. The computational model was an abstraction of reality (Brugiere et al., 2022) and, as such, functioned only as a simplification. The findings were contingent upon synthesized behavioral parameters—like household sensitivity to fines or incentives—drawn from the extant literature (\cite{Chen2023,MuZhang2021}) and triangulated through key informant interviews with barangay officials, not from direct psychometric surveys of individual households. This proxy approach introduces potential divergence between modeled agent behavior and actual resident decision-making. Furthermore, the evaluation of the Deep Reinforcement Learning (DRL) agent's performance is complicated by the absence of a verifiable ground-truth optimal policy; the agent's discovered strategies are assessed relative to heuristic baselines and statistical convergence metrics within the simulation environment, which may not fully translate to real-world governance dynamics. Consequently, the framework was intended strictly as a decision- support tool for exploring policy trade- offs, not as an automated- policy- making mechanism. Furthermore, a major real- world constraint was that the model could not capture the nuanced political or social feasibility of strictly enforcing penalties \citep{Nishimura2022}. This gap between algorithmic optimality and governance reality, which was specifically identified by the MENRO Head (refer to Appendix B.1), confirmed that the DRL- derived optimal policy had to be carefully assessed for its implementation viability before adoption \citep{Dalugdog2021}.